% Options for packages loaded elsewhere
\PassOptionsToPackage{unicode}{hyperref}
\PassOptionsToPackage{hyphens}{url}
\documentclass[
]{article}
\usepackage{xcolor}
\usepackage[margin=2.5cm]{geometry}
\usepackage{amsmath,amssymb}
\setcounter{secnumdepth}{-\maxdimen} % remove section numbering
\usepackage{iftex}
\ifPDFTeX
  \usepackage[T1]{fontenc}
  \usepackage[utf8]{inputenc}
  \usepackage{textcomp} % provide euro and other symbols
\else % if luatex or xetex
  \usepackage{unicode-math} % this also loads fontspec
  \defaultfontfeatures{Scale=MatchLowercase}
  \defaultfontfeatures[\rmfamily]{Ligatures=TeX,Scale=1}
\fi
\usepackage{lmodern}
\ifPDFTeX\else
  % xetex/luatex font selection
\fi
% Use upquote if available, for straight quotes in verbatim environments
\IfFileExists{upquote.sty}{\usepackage{upquote}}{}
\IfFileExists{microtype.sty}{% use microtype if available
  \usepackage[]{microtype}
  \UseMicrotypeSet[protrusion]{basicmath} % disable protrusion for tt fonts
}{}
\makeatletter
\@ifundefined{KOMAClassName}{% if non-KOMA class
  \IfFileExists{parskip.sty}{%
    \usepackage{parskip}
  }{% else
    \setlength{\parindent}{0pt}
    \setlength{\parskip}{6pt plus 2pt minus 1pt}}
}{% if KOMA class
  \KOMAoptions{parskip=half}}
\makeatother
\usepackage{longtable,booktabs,array}
\usepackage{calc} % for calculating minipage widths
% Correct order of tables after \paragraph or \subparagraph
\usepackage{etoolbox}
\makeatletter
\patchcmd\longtable{\par}{\if@noskipsec\mbox{}\fi\par}{}{}
\makeatother
% Allow footnotes in longtable head/foot
\IfFileExists{footnotehyper.sty}{\usepackage{footnotehyper}}{\usepackage{footnote}}
\makesavenoteenv{longtable}
\ifLuaTeX
\usepackage[bidi=basic]{babel}
\else
\usepackage[bidi=default]{babel}
\fi
\babelprovide[main,import]{spanish}
% get rid of language-specific shorthands (see #6817):
\let\LanguageShortHands\languageshorthands
\def\languageshorthands#1{}
\setlength{\emergencystretch}{3em} % prevent overfull lines
\providecommand{\tightlist}{%
  \setlength{\itemsep}{0pt}\setlength{\parskip}{0pt}}
\usepackage{bookmark}
\IfFileExists{xurl.sty}{\usepackage{xurl}}{} % add URL line breaks if available
\urlstyle{same}
\hypersetup{
  pdflang={es},
  hidelinks,
  pdfcreator={LaTeX via pandoc}}

\usepackage{newunicodechar}
\newunicodechar{≥}{\ensuremath{\geq}}
\newunicodechar{→}{\ensuremath{\rightarrow}}
\newunicodechar{↔}{\ensuremath{\leftrightarrow}}

\author{}
\date{}

\begin{document}

{
\setcounter{tocdepth}{3}
\tableofcontents
}
\section{Especificación de Requisitos de Software (SRS) ---
BIOPET}\label{especificaciuxf3n-de-requisitos-de-software-srs-biopet}

\textbf{Versión:} v1.0.0 (Entrega Final, PFC Aplicaciones Web 2026-2027)
\textbf{Conforme a:} ISO/IEC/IEEE 29148:2018 (estructura de SRS), INCOSE
Guide to Writing Requirements v4 (calidad de requisitos individuales y
de conjunto), criterios INVEST de Cohn (historias de usuario, aplicados
explícitamente en la sección 3.3), plantilla de Cockburn (casos de uso).

\textbf{Estado de aprobación:} este documento \textbf{no está firmado}
por el docente-director al momento de esta revisión. El PDF
\texttt{SRS-v1.0.0.pdf} se envía para aprobación/firma tras cerrar esta
revisión; el estado de firma se actualizará aquí, con fecha real,
únicamente cuando exista una firma legítima recibida del docente ---
nunca antes, y nunca simulada o copiada de otro documento.

\textbf{Procedencia de este documento:} este archivo se reconstruye a
partir de dos fuentes: (1) el SRS de la Entrega 1A
(\texttt{PFC\_Entrega1A\_BMT.pdf}), que especificaba el backend sobre
ASP.NET Core 8/C\#, y (2) el código realmente implementado en las
Entregas 1B y Tercera Entrega, sobre \textbf{Java 21 + Spring Boot 3.2 +
Spring Security 6 + Spring Data JPA + PostgreSQL 16 + Redis 7}. Ningún
requisito de este documento fue inventado sin sustento: cada uno
proviene del SRS original, del código fuente verificado, o de la Guía
oficial de la Tercera Entrega. El detalle de cada cambio respecto al
documento original está en
\texttt{docs/requisitos/cambios/CAMBIOS-SRS.md}.

\begin{center}\rule{0.5\linewidth}{0.5pt}\end{center}

\subsection{Índice}\label{uxedndice}

\begin{enumerate}
\def\labelenumi{\arabic{enumi}.}
\tightlist
\item
  Introducción
\item
  Descripción global
\item
  Requisitos específicos 3.1. Requisitos funcionales (REQ-F) 3.2.
  Requisitos no funcionales (REQ-NF) 3.3. Criterios INVEST aplicados a
  las historias de usuario
\item
  Trazabilidad (resumen)
\item
  Modelo de datos (referencia)
\item
  Interfaces de usuario (referencia)
\item
  Observaciones e información pendiente
\end{enumerate}

\begin{center}\rule{0.5\linewidth}{0.5pt}\end{center}

\subsection{1. Introducción}\label{introducciuxf3n}

\subsubsection{1.1. Propósito}\label{propuxf3sito}

Este documento especifica los requisitos funcionales y no funcionales
del sistema BIOPET en su estado \textbf{v1.0.0 (Entrega Final)}. Su
propósito es servir como fuente única de verdad para la matriz de
trazabilidad, las pruebas automatizadas y la evidencia empírica exigidas
por el bloque A.3 de la Guía, incorporando además los módulos de la
Unidad IV (citas, consultas, vacunas, gestión administrativa de usuarios
y consulta externa de especies) que no existían en la revisión v0.9.0-rc
de la Tercera Entrega.

\subsubsection{1.2. Alcance}\label{alcance}

BIOPET es un sistema web para centralizar la información clínica y
administrativa de clínicas veterinarias de pequeña y mediana escala. El
alcance completo del producto (definido en la Entrega 1A) contempla
gestión de dueños y mascotas, historial clínico, citas, telemetría IoT,
recomendación clínica asistida, facturación digital y reportes.

\textbf{Alcance implementado y verificado en v1.0.0:} autenticación
completa (registro, login, refresh, logout con revocación), control de
acceso por rol (RBAC), CRUD completo de la entidad Mascota con
verificación de propiedad, y resumen agregado de mascotas por especie
(vía función SQL). Este es el subconjunto Must Have que el equipo se
comprometió a entregar operativo y reproducible desde la Tercera Entrega
(v0.9.0-rc).

\textbf{Ampliado en v1.0.0 (Unidad IV):} gestión administrativa de
usuarios (REQ-F-023), CRUD de vacunas (REQ-F-024) y consulta externa de
información de especies con caché (REQ-F-025), los tres con interfaz de
usuario real para vacunas
(\texttt{frontend/src/app/features/vacunas.component.ts}) y solo a nivel
de API para usuarios/especies externas (sin pantalla propia todavía).
Adicionalmente, el backend incorporó CRUD completo de citas
(\texttt{CitaController}) y de consultas médicas
(\texttt{ConsultaController}), que formalizan parcialmente REQ-F-015 y
REQ-F-013 respectivamente (detalle de alcance real, sin sobre-declarar,
en la sección 3.1).

\textbf{Alcance pendiente, heredado de la Entrega 1A:} vista consolidada
de historial clínico cronológico y calendario interactivo de citas
(ambos con backend real desde v1.0.0, sin la pieza específica pendiente
--- ver ``Ampliado en v1.0.0'' arriba y la sección 3.1), más
prescripción de medicamentos, telemetría IoT, ubicación en mapa,
recomendaciones asistidas, facturación digital y reportes exportables
(estos últimos seis, sin ningún código de backend ni de frontend
todavía). El modelo de datos conceptual de estos módulos ya existe (ver
sección 5 y el DER de la Entrega 1A). Se documentan como requisitos con
estado ``pendiente'' o ``parcial'' (según corresponda) para que la
matriz de trazabilidad (bloque A.3.3 de la Guía) los declare
correctamente, en vez de omitirlos.

\subsubsection{1.3. Definiciones, acrónimos y
abreviaturas}\label{definiciones-acruxf3nimos-y-abreviaturas}

\begin{longtable}[]{@{}
  >{\raggedright\arraybackslash}p{(\linewidth - 2\tabcolsep) * \real{0.5000}}
  >{\raggedright\arraybackslash}p{(\linewidth - 2\tabcolsep) * \real{0.5000}}@{}}
\toprule\noalign{}
\begin{minipage}[b]{\linewidth}\raggedright
Término
\end{minipage} & \begin{minipage}[b]{\linewidth}\raggedright
Definición
\end{minipage} \\
\midrule\noalign{}
\endhead
\bottomrule\noalign{}
\endlastfoot
RBAC & Role-Based Access Control: control de acceso basado en el rol del
usuario autenticado. \\
JWT & JSON Web Token (RFC 7519): token firmado que representa
afirmaciones sobre un sujeto autenticado. \\
ORM & Object-Relational Mapping: mapeo objeto-relacional (Spring Data
JPA / Hibernate en este proyecto). \\
SP & Stored Procedure / función almacenada en el motor de base de
datos. \\
TTL & Time To Live: tiempo de vida de una entrada de caché. \\
MoSCoW & Técnica de priorización: Must, Should, Could, Won't. \\
shall & Verbo modal normativo de ISO/IEC/IEEE 29148 que indica
obligación contractual del requisito. \\
\end{longtable}

\subsubsection{1.4. Referencias}\label{referencias}

\begin{itemize}
\tightlist
\item
  ISO/IEC/IEEE 29148:2018 --- Requirements Engineering.
\item
  INCOSE Guide to Writing Requirements v4.
\item
  RFC 7519 (JWT), RFC 7807 (ProblemDetails).
\item
  Guía de la Tercera Entrega --- PFC Aplicaciones Web 2026-2027, UTEQ.
\item
  \texttt{PFC\_Entrega1A\_BMT.pdf} --- SRS y diseño original (ASP.NET
  Core).
\item
  \texttt{docs/requisitos/cambios/CAMBIOS-SRS.md} --- registro de
  cambios respecto al documento original.
\end{itemize}

\subsubsection{1.5. Resumen del documento}\label{resumen-del-documento}

La sección 2 describe el producto de forma global (perspectiva,
funciones, usuarios, restricciones). La sección 3 contiene el detalle de
cada requisito funcional y no funcional con su identificador
persistente, prioridad MoSCoW, criterio de aceptación y método de
verificación. La sección 4 resume la trazabilidad end-to-end. Las
secciones 5 y 6 remiten a los artefactos de modelo de datos e interfaz
que ya existen en el repositorio. La sección 7 documenta, sin inventar
contenido, lo que aún falta por completar.

\begin{center}\rule{0.5\linewidth}{0.5pt}\end{center}

\subsection{2. Descripción global}\label{descripciuxf3n-global}

\subsubsection{2.1. Perspectiva del
producto}\label{perspectiva-del-producto}

BIOPET es un sistema nuevo, no una extensión de un producto previo
existente en la clínica. Reemplaza procesos manuales (hojas de cálculo,
mensajería) por una plataforma web centralizada, según el problema
documentado en la Entrega 1A (entrevistas a tres veterinarios y dos
auxiliares, encuestas a quince dueños de mascotas).

\subsubsection{2.2. Cambio de plataforma tecnológica respecto a la
Entrega
1A}\label{cambio-de-plataforma-tecnoluxf3gica-respecto-a-la-entrega-1a}

El SRS original (Entrega 1A) especificaba ASP.NET Core 8/C\# como
backend (ADR-001 original) y Bootstrap 5 + HTML/CSS/JS como frontend. El
equipo migró la implementación real a \textbf{Java 21 + Spring Boot 3.2}
en el backend y \textbf{Angular 17+} en el frontend, decisión reflejada
en \texttt{ADR-002-pila-tecnologica.md} del repositorio actual. Los
requisitos funcionales de negocio (qué hace el sistema) no cambiaron;
los requisitos no funcionales técnicos (cómo se implementa: framework,
ORM, mecanismo JWT) se actualizaron para reflejar la pila real, evitando
que el SRS describa un sistema que no es el que existe en el
repositorio.

\subsubsection{2.3. Funciones del producto
(resumen)}\label{funciones-del-producto-resumen}

\begin{itemize}
\tightlist
\item
  Gestión de usuarios y autenticación (roles: \texttt{ADMIN},
  \texttt{VETERINARIO}, \texttt{AUXILIAR}, \texttt{DUENO}), incluida la
  administración de cuentas por un \texttt{ADMIN} (Unidad IV).
\item
  Gestión de mascotas (CRUD + resumen agregado por especie).
\item
  Gestión de vacunas (CRUD, con interfaz de usuario) y de citas y
  consultas médicas (CRUD a nivel de API; sin calendario interactivo ni
  vista de historial clínico consolidado todavía --- ver 3.1).
\item
  Consulta de información externa de especies (taxonomía, hábitat,
  dieta), con caché.
\item
  \emph{(Pendiente)} Vista consolidada de historial clínico cronológico,
  prescripción de medicamentos, telemetría IoT, recomendaciones
  asistidas, facturación y reportes exportables.
\end{itemize}

\subsubsection{2.4. Características de los
usuarios}\label{caracteruxedsticas-de-los-usuarios}

\begin{longtable}[]{@{}
  >{\raggedright\arraybackslash}p{(\linewidth - 4\tabcolsep) * \real{0.3333}}
  >{\raggedright\arraybackslash}p{(\linewidth - 4\tabcolsep) * \real{0.3333}}
  >{\raggedright\arraybackslash}p{(\linewidth - 4\tabcolsep) * \real{0.3333}}@{}}
\toprule\noalign{}
\begin{minipage}[b]{\linewidth}\raggedright
Rol
\end{minipage} & \begin{minipage}[b]{\linewidth}\raggedright
Descripción
\end{minipage} & \begin{minipage}[b]{\linewidth}\raggedright
Nivel técnico esperado
\end{minipage} \\
\midrule\noalign{}
\endhead
\bottomrule\noalign{}
\endlastfoot
\texttt{ADMIN} & Supervisa roles, operatividad general del sistema y
tiene visibilidad global de datos. & Medio-alto (personal administrativo
de la clínica). \\
\texttt{VETERINARIO} & Registra diagnósticos y gestiona historiales
clínicos (módulo pendiente); gestiona mascotas. & Medio (personal
clínico). \\
\texttt{AUXILIAR} & Apoya operaciones administrativas y de gestión de
mascotas. & Medio. \\
\texttt{DUENO} & Consulta únicamente la información de sus propias
mascotas. & Básico (público general, sin capacitación previa). \\
\end{longtable}

\subsubsection{2.5. Restricciones}\label{restricciones}

\begin{itemize}
\tightlist
\item
  El proyecto debe completarse dentro de un semestre académico (PPA
  2026-2027), lo que limita el alcance implementado al núcleo Must Have.
\item
  La base de datos debe ser PostgreSQL 16 (decisión ya tomada, ver
  \texttt{ADR-004-postgresql.md}).
\item
  La autenticación debe usar JWT en cookies
  \texttt{HttpOnly\ +\ Secure\ +\ \ \ SameSite=Strict} (bloque A.1 de la
  Guía), no \texttt{localStorage} ni \texttt{sessionStorage}.
\item
  Toda operación de base de datos que no sea CRUD elemental debe
  implementarse como función/procedimiento almacenado (bloque A.2 de la
  Guía), no como JPQL/HQL con joins o agregaciones.
\end{itemize}

\subsubsection{2.6. Supuestos y
dependencias}\label{supuestos-y-dependencias}

\begin{itemize}
\tightlist
\item
  Se asume disponibilidad de un motor PostgreSQL 16 y Redis 7 accesibles
  desde el backend (vía Docker Compose en desarrollo).
\item
  Se asume que los servicios externos mencionados en la Entrega 1A (IoT,
  IA cognitiva, correo) se integrarán en fases posteriores; ninguno está
  implementado todavía y no se documentan requisitos técnicos de
  integración hasta que exista una decisión de arquitectura formal para
  ellos.
\end{itemize}

\begin{center}\rule{0.5\linewidth}{0.5pt}\end{center}

\subsection{3. Requisitos específicos}\label{requisitos-especuxedficos}

Cada requisito seguirá el patrón
\texttt{{[}condición{]}\ {[}sujeto{]}\ shall\ {[}acción{]}\ {[}objeto{]}\ {[}restricción{]}}
de ISO/IEC/IEEE 29148, con identificador persistente, rationale,
prioridad MoSCoW, criterio de aceptación medible y método de
verificación, cumpliendo las características INCOSE C1--C9 (Necessary,
Appropriate, Unambiguous, Complete, Singular, Feasible, Verifiable,
Correct, Conforming).

\subsubsection{3.1. Requisitos funcionales
(REQ-F)}\label{requisitos-funcionales-req-f}

\begin{quote}
\textbf{Nota de consistencia (corrección respecto a una versión anterior
de este documento):} la numeración de esta sección es la que ya está
fijada en \texttt{docs/requisitos/historias/HistoriasUsuario.md} (HU-001
a HU-020) y \texttt{docs/requisitos/casos-de-uso/CasosDeUso.md} (CU-01 a
CU-20), documentos que ya existían en el repositorio con 20 requisitos
funcionales completos (incluyendo uno propio para ``consultar mascota
por id'' y otro para ``auditoría'', que una reconstrucción anterior de
este SRS no tenía desglosados igual). Esta versión del SRS adopta esa
numeración como única fuente de verdad, para que los tres documentos no
se contradigan entre sí.
\end{quote}

\paragraph{Implementados y verificados (REQ-F-001 a REQ-F-012,
REQ-F-021)}\label{implementados-y-verificados-req-f-001-a-req-f-012-req-f-021}

\textbf{REQ-F-001 --- Registro de usuario dueño de mascota} -
\textbf{Tipo:} Funcional · \textbf{Prioridad:} Must -
\textbf{Enunciado:} El sistema deberá permitir que cualquier visitante
se registre proporcionando nombre, correo electrónico y contraseña,
asignándole automáticamente el rol \texttt{ROLE\_DUENO}, ignorando
cualquier rol enviado por el cliente. - \textbf{Rationale:} heredado de
RF-01 de la Entrega 1A; el registro público solo aplica al rol dueño,
los demás roles se crean administrativamente. - \textbf{Verificación:}
Test \texttt{AuthControllerTest} (registro exitoso). -
\textbf{Trazabilidad:} → HU-001 → CU-01 →
\texttt{AuthController.registro} → \texttt{AuthService.registrar} →
\texttt{POST\ /api/auth/registro}. - \textbf{Estado:} verificado.

\textbf{REQ-F-002 --- Rechazo de registro con correo duplicado} -
\textbf{Tipo:} Funcional · \textbf{Prioridad:} Must -
\textbf{Enunciado:} Al recibir una solicitud de registro con un correo
electrónico ya existente, el sistema deberá rechazarla con un código
\texttt{409\ Conflict} y un cuerpo \texttt{ProblemDetails}, sin crear
ningún registro. - \textbf{Rationale:} funcionalidad presente en el
código (\texttt{EmailDuplicadoException}) no documentada como requisito
independiente en el SRS original de la Entrega 1A; se agrega para cerrar
el hueco (bloque A.3 de la Guía). - \textbf{Verificación:} Test
\texttt{AuthControllerTest} (registro con correo duplicado). -
\textbf{Trazabilidad:} → HU-001 → CU-01 → \texttt{AuthService.registrar}
→ \texttt{EmailDuplicadoException} → \texttt{GlobalExceptionHandler}. -
\textbf{Estado:} verificado.

\textbf{REQ-F-003 --- Autenticación mediante usuario y contraseña} -
\textbf{Tipo:} Funcional · \textbf{Prioridad:} Must -
\textbf{Enunciado:} El sistema deberá permitir que un usuario registrado
inicie sesión mediante correo electrónico y contraseña, emitiendo un
access token (1 hora de vigencia) y un refresh token con los siete
claims JWT estándar. - \textbf{Rationale:} heredado de RF-16/RF-WEB-01
de la Entrega 1A. - \textbf{Verificación:} Test
\texttt{AuthControllerTest}, \texttt{JwtServiceTest}. -
\textbf{Trazabilidad:} → HU-002 → CU-02 → \texttt{AuthController.login}
→ \texttt{POST\ /api/auth/login}. - \textbf{Estado:} verificado.

\textbf{REQ-F-004 --- Renovación de sesión (refresh)} - \textbf{Tipo:}
Funcional · \textbf{Prioridad:} Must - \textbf{Enunciado:} El sistema
deberá permitir renovar el access token utilizando un refresh token
válido y no revocado, sin exigir nuevamente las credenciales del
usuario. - \textbf{Rationale:} heredado de RNF-03/RNF-WEB-03 de la
Entrega 1A (expiración configurable de JWT), llevado a requisito
funcional explícito. - \textbf{Verificación:} Test
\texttt{AuthControllerTest} (flujo refresh). - \textbf{Trazabilidad:} →
HU-003 → CU-03 → \texttt{AuthController.refresh} →
\texttt{POST\ /api/auth/refresh}. - \textbf{Estado:} verificado.

\textbf{REQ-F-005 --- Cierre de sesión con revocación de token} -
\textbf{Tipo:} Funcional · \textbf{Prioridad:} Must -
\textbf{Enunciado:} El sistema deberá permitir cerrar sesión,
registrando el \texttt{jti} del token vigente en una lista negra en
Redis con TTL igual a su tiempo restante de expiración, de modo que una
solicitud posterior con ese token a un recurso protegido responda
\texttt{401}. - \textbf{Rationale:} heredado de RF-17/RF-WEB-04 de la
Entrega 1A; mecanismo de revocación definido en
\texttt{ADR-003-jwt-redis.md}. - \textbf{Verificación:} Test
\texttt{AuthControllerTest}, \texttt{AuthenticationAuditServiceTest}. -
\textbf{Trazabilidad:} → HU-004 → CU-04 → \texttt{AuthController.logout}
→ \texttt{TokenBlacklistService.revoke} →
\texttt{POST\ /api/auth/logout}. - \textbf{Estado:} verificado.

\textbf{REQ-F-006 --- Control de acceso por rol (RBAC)} - \textbf{Tipo:}
Funcional · \textbf{Prioridad:} Must - \textbf{Enunciado:} El sistema
deberá restringir el acceso a cada endpoint protegido según el rol del
usuario autenticado, rechazando con \texttt{403\ Forbidden} cualquier
solicitud de un rol no autorizado para ese recurso. -
\textbf{Rationale:} heredado de RF-13/RF-WEB-02 de la Entrega 1A; caso
de uso transversal incluido implícitamente por la gestión de mascotas. -
\textbf{Verificación:} Test \texttt{MascotaControllerTest} (casos 403
por rol). - \textbf{Trazabilidad:} → HU-005 → CU-05 → anotaciones
\texttt{@PreAuthorize} en \texttt{MascotaController}. - \textbf{Estado:}
verificado.

\textbf{REQ-F-007 --- Consulta del perfil propio} - \textbf{Tipo:}
Funcional · \textbf{Prioridad:} Should - \textbf{Enunciado:} El sistema
deberá permitir que un usuario autenticado consulte sus propios datos de
perfil (nombre, correo, rol), sin exponer el listado completo de
usuarios. - \textbf{Rationale:} funcionalidad presente en el código
(\texttt{GET\ /api/usuarios/me}), base del \texttt{authGuard} de
Angular; no documentada como requisito independiente en el SRS original.
- \textbf{Verificación:} inspección de \texttt{UsuarioController} y
prueba manual con Postman (sin test automatizado formal registrado). -
\textbf{Trazabilidad:} → HU-006 → CU-06 → \texttt{UsuarioController.me}
→ \texttt{AuthService.perfil}. - \textbf{Estado:} verificado.

\textbf{REQ-F-008 --- Creación de mascota} - \textbf{Tipo:} Funcional ·
\textbf{Prioridad:} Must - \textbf{Enunciado:} El sistema deberá
permitir que un usuario con rol \texttt{ADMIN}, \texttt{VETERINARIO} o
\texttt{AUXILIAR} registre una nueva mascota (nombre, especie, raza,
fecha de nacimiento) asociada a un dueño existente con rol
\texttt{ROLE\_DUENO}. - \textbf{Rationale:} heredado de RF-01/RF-02 de
la Entrega 1A (registro y asociación con propietario), acotado a la
entidad Mascota implementada en esta fase. - \textbf{Verificación:} Test
\texttt{MascotaControllerTest} (creación exitosa y validación de
campos). - \textbf{Trazabilidad:} → HU-007 → CU-07 →
\texttt{MascotaController.crear} → \texttt{POST\ /api/mascotas}. -
\textbf{Estado:} verificado.

\textbf{REQ-F-009 --- Listado paginado de mascotas activas, con filtrado
por propietario según rol} - \textbf{Tipo:} Funcional ·
\textbf{Prioridad:} Must - \textbf{Enunciado:} El sistema deberá
permitir consultar el listado paginado de mascotas activas; si el
solicitante tiene rol \texttt{ROLE\_DUENO}, el listado deberá
restringirse únicamente a sus propias mascotas; para los demás roles
autorizados, el listado deberá incluir todas las mascotas activas. -
\textbf{Rationale:} heredado de RF-02 de la Entrega 1A, con la regla de
aislamiento de datos entre dueños explícita en el código
(\texttt{MascotaService.listar}). - \textbf{Verificación:} Test
\texttt{MascotaControllerTest} (listado por rol \texttt{ROLE\_DUENO}
vs.~resto). - \textbf{Trazabilidad:} → HU-008 → CU-08 →
\texttt{MascotaController.listar} → \texttt{GET\ /api/mascotas}. -
\textbf{Estado:} verificado.

\textbf{REQ-F-010 --- Consulta de mascota por identificador} -
\textbf{Tipo:} Funcional · \textbf{Prioridad:} Must -
\textbf{Enunciado:} El sistema deberá permitir consultar el detalle
completo de una mascota a partir de su identificador, respondiendo
\texttt{404} con \texttt{ProblemDetails} si no existe. -
\textbf{Rationale:} complementa REQ-F-009 para el caso de consulta
puntual, necesario antes de una edición o de mostrar el detalle en la
interfaz; no estaba desglosado como requisito independiente en una
versión anterior de este documento. - \textbf{Verificación:} Test
\texttt{MascotaControllerTest} y \texttt{RecursoNoEncontradoException}.
- \textbf{Trazabilidad:} → HU-009 → CU-09 →
\texttt{MascotaController.buscar} → \texttt{GET\ /api/mascotas/\{id\}}.
- \textbf{Estado:} verificado.

\textbf{REQ-F-011 --- Actualización de mascota con verificación de
propiedad} - \textbf{Tipo:} Funcional · \textbf{Prioridad:} Must -
\textbf{Enunciado:} El sistema deberá permitir actualizar los datos de
una mascota existente a un usuario con rol \texttt{ADMIN},
\texttt{VETERINARIO} o \texttt{AUXILIAR}, registrando la fecha de
actualización (\texttt{actualizado\_en}). - \textbf{Rationale:} heredado
de RF-02 de la Entrega 1A. - \textbf{Verificación:} Test
\texttt{MascotaControllerTest} (actualización exitosa y caso 403/404). -
\textbf{Trazabilidad:} → HU-010 → CU-10 →
\texttt{MascotaController.actualizar} →
\texttt{PUT\ /api/mascotas/\{id\}}. - \textbf{Estado:} verificado.

\textbf{REQ-F-012 --- Baja lógica de una mascota} - \textbf{Tipo:}
Funcional · \textbf{Prioridad:} Must - \textbf{Enunciado:} El sistema
deberá permitir dar de baja lógicamente (cambio del atributo
\texttt{activo} a \texttt{false}) una mascota existente, sin eliminar
físicamente el registro de la base de datos. - \textbf{Rationale:}
preservación de historial e integridad referencial futura con historial
clínico y citas (módulos pendientes). - \textbf{Verificación:} Test
\texttt{MascotaControllerTest} (baja y verificación de que la mascota
deja de aparecer en el listado activo). - \textbf{Trazabilidad:} →
HU-011 → CU-11 → \texttt{MascotaController.eliminar} →
\texttt{DELETE\ /api/mascotas/\{id\}}. - \textbf{Estado:} verificado.

\textbf{REQ-F-021 --- Resumen de mascotas activas agrupadas por especie}
- \textbf{Tipo:} Funcional · \textbf{Prioridad:} Should -
\textbf{Enunciado:} El sistema deberá permitir consultar el total de
mascotas activas agrupadas por especie; para roles distintos de
\texttt{ADMIN}, el resultado deberá restringirse siempre a las mascotas
del propio usuario autenticado; para el rol \texttt{ADMIN}, el resultado
deberá poder filtrarse por un dueño específico o consultarse sin filtro
para obtener el total global. - \textbf{Rationale:} nuevo requisito de
la Tercera Entrega, exigido por el bloque A.2.2 de la Guía como ejemplo
de operación agregada (\texttt{GROUP\ BY}) que debe implementarse
obligatoriamente vía función/procedimiento almacenado, no vía JPQL. Se
numera 021 (no 013) porque los identificadores 013 a 020 ya estaban
ocupados por los requisitos heredados de la Entrega 1A, fijados en
\texttt{HistoriasUsuario.md}/\texttt{CasosDeUso.md} antes de que este
requisito existiera. - \textbf{Verificación:} Test de integración
\texttt{ResumenEspeciesIntegrationTest} (Testcontainers, PostgreSQL
real). - \textbf{Trazabilidad:} → HU-020 → CU-20 →
\texttt{MascotaController.resumenPorEspecies} →
\texttt{MascotaService.resumenPorEspecie} →
\texttt{MascotaRepository.resumenPorEspecie} → función
\texttt{fn\_resumen\_mascotas\_por\_especie} →
\texttt{GET\ /api/mascotas/resumen-especies}. - \textbf{Tipo de acceso a
datos:} SP (agregación con \texttt{GROUP\ BY}, prohibido en ORM
elemental por el bloque A.2.1 de la Guía). - \textbf{Estado:}
verificado.

\paragraph{Pendientes, heredados de la Entrega 1A (REQ-F-013 a
REQ-F-020)}\label{pendientes-heredados-de-la-entrega-1a-req-f-013-a-req-f-020}

\begin{quote}
Estos requisitos provienen directamente de RF-03 a RF-15 del SRS
original y están desglosados exactamente como en
\texttt{HistoriasUsuario.md} (HU-012 a HU-019) y \texttt{CasosDeUso.md}
(CU-12 a CU-19). El modelo de datos conceptual ya existe en el DER de la
Entrega 1A (entidades \texttt{Cita}, \texttt{Historial\_Clinico},
\texttt{Dispositivo\_IoT}, \texttt{Chat\_Triage},
\texttt{Producto\_Servicio}, \texttt{Factura}, etc.), pero \textbf{no
hay código de backend ni de frontend implementado} para ninguno de ellos
en v0.9.0-rc. Se reclasifican de prioridad Must/Media (documento
original) a Should/Could, porque no forman parte del compromiso de esta
entrega, y se marcan como ``pendiente'' para que la matriz de
trazabilidad los declare sin ocultarlos.
\end{quote}

\begin{longtable}[]{@{}
  >{\raggedright\arraybackslash}p{(\linewidth - 10\tabcolsep) * \real{0.1667}}
  >{\raggedright\arraybackslash}p{(\linewidth - 10\tabcolsep) * \real{0.1667}}
  >{\raggedright\arraybackslash}p{(\linewidth - 10\tabcolsep) * \real{0.1667}}
  >{\raggedright\arraybackslash}p{(\linewidth - 10\tabcolsep) * \real{0.1667}}
  >{\raggedright\arraybackslash}p{(\linewidth - 10\tabcolsep) * \real{0.1667}}
  >{\raggedright\arraybackslash}p{(\linewidth - 10\tabcolsep) * \real{0.1667}}@{}}
\toprule\noalign{}
\begin{minipage}[b]{\linewidth}\raggedright
Id
\end{minipage} & \begin{minipage}[b]{\linewidth}\raggedright
Descripción (resumen)
\end{minipage} & \begin{minipage}[b]{\linewidth}\raggedright
Prioridad
\end{minipage} & \begin{minipage}[b]{\linewidth}\raggedright
Origen
\end{minipage} & \begin{minipage}[b]{\linewidth}\raggedright
HU / CU
\end{minipage} & \begin{minipage}[b]{\linewidth}\raggedright
Estado
\end{minipage} \\
\midrule\noalign{}
\endhead
\bottomrule\noalign{}
\endlastfoot
REQ-F-013 & Registrar atención médica y almacenar/consultar el historial
clínico de cada mascota de forma cronológica. & Should & RF-03, RF-04 &
HU-012 / CU-12 & parcial \\
REQ-F-014 & Registrar medicamentos prescritos durante una atención
médica. & Could & RF-05 & HU-013 / CU-13 & pendiente \\
REQ-F-015 & Registrar, modificar y consultar citas veterinarias mediante
calendario interactivo. & Should & RF-06 & HU-014 / CU-14 & parcial \\
REQ-F-016 & Proveer una API para recibir datos de dispositivos IoT de
rastreo asociados a mascotas. & Could & RF-08, RF-09 & HU-015 / CU-15 &
pendiente \\
REQ-F-017 & Generar recomendaciones clínicas informativas a partir del
historial médico (enunciado ampliado --- ver bloque detallado, cierre de
OBS-04). & Could & RF-10 & HU-016 / CU-16 & pendiente \\
REQ-F-018 & Generar comprobantes de pago digitales en PDF y registrar el
método de pago de un servicio veterinario. & Should & RF-11, RF-12 &
HU-017 / CU-17 & pendiente \\
REQ-F-019 & Generar reportes estadísticos exportables en PDF y Excel. &
Could & RF-14 & HU-018 / CU-18 & pendiente \\
REQ-F-020 & Registrar las operaciones relevantes del sistema (usuario,
fecha, hora) para fines de auditoría. & Should & RF-15 & HU-019 / CU-19
& parcial \\
REQ-F-022 & Enviar notificaciones al usuario por correo electrónico ante
eventos relevantes de su cuenta. & Could & RF-07 & HU-021 / CU-21 &
pendiente \\
\end{longtable}

\begin{quote}
\textbf{Nota sobre REQ-F-013 y REQ-F-015 (revisión v1.0.0 --- código
actual):} a diferencia de los demás requisitos de esta tabla, no se
dejan como ``pendiente'' en v1.0.0 porque la Unidad IV sí incorporó
código real para ambos, verificado contra el repositorio actual (paso de
revisión obligatorio antes de cerrar esta entrega):
\texttt{CitaController}/ \texttt{CitaService}/\texttt{CitaRepository}
(CRUD completo: listar, buscar, crear, actualizar, eliminar) y
\texttt{ConsultaController}/\texttt{ConsultaService}/
\texttt{ConsultaRepository} (mismo CRUD), ambos con clase de test
dedicada (\texttt{CitaControllerTest}, \texttt{ConsultaControllerTest}).
Se marcan \textbf{``parcial''}, no ``verificado'', porque el enunciado
original exige más de lo que el CRUD por sí solo cubre: REQ-F-015 pide
explícitamente un ``calendario interactivo'' (no existe pantalla de
citas en \texttt{frontend/src/app/features/}, solo la API) y REQ-F-013
pide el historial clínico ``de forma cronológica'' (existe la proyección
\texttt{HistorialClinico}/función
\texttt{fn\_historial\_clinico\_mascota} a nivel de repositorio, pero
ningún controlador la expone todavía como endpoint consolidado --- la
única vía real hoy es el CRUD genérico de consultas, sin agrupar ni
ordenar por mascota). \texttt{docs/trazabilidad/matriz.csv} registra
además, para REQ-F-013, que \texttt{ConsultaService.listar()} \textbf{no
filtra por propietario para \texttt{ROLE\_DUENO}} pese a la regla
general de aislamiento de datos que sí aplican
\texttt{MascotaService}/\texttt{VacunaService} (ver REQ-F-009); y que
\texttt{GET\ /api/consultas} (listado) y
\texttt{PUT\ /api/consultas/\{id\}} no tienen prueba automatizada
dedicada en \texttt{ConsultaControllerTest}. Ambos quedan como
limitación explícita en la sección 7, no como pendientes ocultos ni como
``verificado'' sin serlo del todo.
\end{quote}

\begin{quote}
\textbf{Nota sobre REQ-F-020 (auditoría):} a diferencia de los demás
requisitos de esta tabla, no está 100\% pendiente: el registro de
eventos de \emph{autenticación} (login exitoso/fallido, bloqueo por rate
limit, refresh, logout) ya opera en producción y está cubierto por
REQ-NF-009 (\texttt{AuthenticationAuditService}). Lo que falta es la
auditoría genérica de operaciones CRUD sobre mascotas y los módulos
pendientes. Se marca ``parcial'', no ``pendiente'', para no ocultar el
trabajo ya hecho.
\end{quote}

\begin{quote}
\textbf{Nota sobre REQ-F-017 (cierre de OBS-04):} la retroalimentación
oficial del SGA de la Entrega 1A señaló, como observación leve, la
ambigüedad de la redacción original de RF-10 (``recomendaciones
informativas''; ver \texttt{docs/observaciones/OBSERVACIONES.md},
OBS-04). El identificador y el alcance no cambian (sigue siendo
\texttt{REQ-F-017}, prioridad \texttt{Could}, estado \texttt{pendiente},
dependiente de \texttt{REQ-F-013}); lo que cambia es exclusivamente la
redacción, reformulada en patrón ``El sistema deberá\ldots{}'' con
entradas, resultado esperado y criterios de aceptación verificables, sin
ampliar la funcionalidad prevista originalmente. El detalle completo
reemplaza al resumen de una sola línea que tenía esta fila en versiones
anteriores del SRS.
\end{quote}

\textbf{REQ-F-017 --- Generación de recomendaciones clínicas
informativas a partir del historial médico} - \textbf{Tipo:} Funcional ·
\textbf{Prioridad:} Could - \textbf{Enunciado:} Al recibir una solicitud
de recomendaciones para una mascota con historial clínico registrado, el
sistema deberá generar, mediante un servicio de IA externo, una lista de
recomendaciones de cuidado en texto a partir de los datos del historial
clínico de esa mascota, y deberá devolver cada recomendación acompañada
de la advertencia explícita ``informativa, no sustituye diagnóstico
veterinario''. - \textbf{Entradas:} identificador de una mascota con
historial clínico registrado (depende de \texttt{REQ-F-013}, módulo
pendiente). - \textbf{Resultado esperado:} una lista de cero a N
recomendaciones en texto plano, cada una acompañada de la advertencia de
carácter informativo; si la mascota no tiene historial clínico
registrado, el sistema deberá responder con una lista vacía, no con un
error. - \textbf{Criterios de aceptación (verificables):} 1. Dado un
historial clínico existente para la mascota solicitada, cuando se invoca
el endpoint de recomendaciones, entonces la respuesta incluye el campo
\texttt{recomendaciones:\ string{[}{]}} y, por cada elemento, el campo
\texttt{advertencia:\ "informativa,\ no\ sustituye\ diagnóstico\ veterinario"}.
2. Dado que la mascota no tiene historial clínico registrado, cuando se
solicita el endpoint de recomendaciones, entonces la respuesta es
\texttt{200\ OK} con \texttt{recomendaciones:\ {[}{]}}. 3. El proveedor
concreto del servicio de IA es una decisión de arquitectura pendiente;
no se implementa en v0.9.0-rc. - \textbf{Rationale:} heredado de RF-10
de la Entrega 1A; reformulado en la Tercera Entrega para cerrar OBS-04
(ambigüedad leve señalada por el docente en ``recomendaciones
informativas''), sin ampliar el alcance original: sigue dependiendo de
\texttt{REQ-F-013} (historial clínico, módulo pendiente) y del mismo
servicio de IA externo ya previsto desde la Entrega 1A. -
\textbf{Verificación:} pendiente --- no implementado en v0.9.0-rc
(depende de \texttt{REQ-F-013}). Método de verificación previsto: test
de integración con un doble de prueba (mock) del servicio de IA,
validando el contrato de entrada/salida y el caso sin historial clínico.
- \textbf{Trazabilidad:} → HU-016 → CU-16 → (módulo pendiente; depende
de \texttt{REQ-F-013}). - \textbf{Estado:} pendiente (redacción cerrada;
implementación no iniciada).

\begin{quote}
\textbf{Nota sobre REQ-F-022 (cierre de OBS-02):} este requisito no
forma parte de la secuencia original RF-03 a RF-15; se numera 022
(siguiente identificador libre después de REQ-F-021) porque corresponde
a \textbf{RF-07}, el requisito que la retroalimentación oficial del SGA
de la Entrega 1A señaló como ausente de la lista consolidada (``FALTA
RF-07 en la lista consolidada, salta RF-06 -\textgreater{} RF-08''; ver
\texttt{docs/observaciones/OBSERVACIONES.md}, OBS-02). RF-07 no estaba
perdido: el diagrama de contexto C4 Nivel 1
(\texttt{docs/diagrams/c4-contexto/C4-L1-contexto.md}, sección
``Trazabilidad'') ya lo identificaba como el ``Servicio de Correos''
mencionado en la Entrega 1A, sin que ninguna versión anterior de este
SRS le hubiera asignado un identificador \texttt{REQ-F} propio. Se le
asigna aquí \texttt{REQ-F-022} (no \texttt{REQ-F-007}, que ya está
ocupado por un requisito distinto --- ``Consulta del perfil propio'' ---
para no duplicar identificadores) con su enunciado completo, criterios
de aceptación verificables y trazabilidad propia (ver bloque detallado
más abajo).
\end{quote}

\textbf{REQ-F-022 --- Notificaciones al usuario por correo electrónico
(Servicio de Correos, RF-07 recuperado)} - \textbf{Tipo:} Funcional ·
\textbf{Prioridad:} Could - \textbf{Enunciado:} Al ocurrir un evento
relevante para la cuenta de un usuario (por ejemplo, registro exitoso),
el sistema deberá enviar una notificación por correo electrónico al
usuario afectado a través de un servicio de correo externo, de forma
asíncrona, sin bloquear ni condicionar la respuesta HTTP de la operación
que originó el evento. - \textbf{Entradas:} un evento de dominio
notificable (tipo de evento, correo electrónico destinatario, datos
contextuales mínimos del evento). - \textbf{Resultado esperado:} el
correo se envía o se encola para envío de forma asíncrona; un fallo del
servicio de correo externo no debe revertir ni hacer fallar la operación
de dominio que originó el evento. - \textbf{Criterios de aceptación
(verificables):} 1. Dado un registro de usuario exitoso, cuando el
servicio de correo está disponible, entonces se envía un correo de
bienvenida a la dirección registrada. 2. Dado que el servicio de correo
externo no está disponible, cuando ocurre un evento notificable,
entonces la operación de dominio que lo originó (por ejemplo, el
registro) completa exitosamente y el fallo de envío queda registrado en
el log del sistema, sin propagarse como error al cliente. 3. El
proveedor de correo concreto (SMTP propio, SES, SendGrid u otro) es una
decisión de arquitectura pendiente; no está implementado en v0.9.0-rc. -
\textbf{Rationale:} corresponde al RF-07 original de la Entrega 1A
(``Servicio de Correos''), documentado como sistema externo en
\texttt{docs/diagrams/c4-contexto/C4-L1-contexto.md} pero sin requisito
\texttt{REQ-F} formal hasta esta revisión. Se incorpora explícitamente
para cerrar OBS-02 sin duplicar \texttt{REQ-F-007} (funcionalidad
distinta ya implementada). - \textbf{Verificación:} pendiente --- no
implementado en v0.9.0-rc. Método de verificación previsto: test de
integración con un servidor SMTP de pruebas (por ejemplo, Mailhog) que
confirme el envío en el caso exitoso y el comportamiento no bloqueante
ante fallo del servicio externo. - \textbf{Trazabilidad:} → HU-021 →
CU-21 → (módulo pendiente; sin controlador ni servicio implementados en
el backend). - \textbf{Estado:} pendiente.

\paragraph{Incorporados en Unidad IV (REQ-F-023 a
REQ-F-025)}\label{incorporados-en-unidad-iv-req-f-023-a-req-f-025}

\begin{quote}
\textbf{Nota de alcance:} los tres requisitos siguientes formalizan
funcionalidades reales incorporadas al backend durante la Unidad IV
(Citas y Consultas ya tenían requisito formal propio ---
\texttt{REQ-F-015} y \texttt{REQ-F-013} respectivamente --- y no se
duplican aquí). Cada uno agrupa el módulo completo, no una operación
HTTP individual: la trazabilidad hacia cada endpoint real se detalla en
\texttt{docs/trazabilidad/matriz.csv}.
\end{quote}

\textbf{REQ-F-023 --- Gestión administrativa de usuarios} -
\textbf{Tipo:} Funcional · \textbf{Prioridad:} Should -
\textbf{Enunciado:} El sistema deberá permitir que un usuario con rol
\texttt{ADMIN} liste, consulte por identificador, cree, actualice y dé
de baja lógica cuentas de usuario de cualquier rol, sin permitir que un
administrador modifique el rol de su propia cuenta. -
\textbf{Rationale:} nuevo requisito de la actividad de Unidad IV;
formaliza el CRUD administrativo de
\texttt{UsuarioController}/\texttt{UsuarioService}, funcionalidad
distinta de \texttt{REQ-F-007} (consulta del perfil propio,
\texttt{/api/usuarios/me}), que no se modifica ni se duplica. Se numera
023 (siguiente identificador libre después de \texttt{REQ-F-022}). -
\textbf{Verificación:} Test \texttt{UsuarioControllerTest} (16 pruebas:
listado, consulta, creación, actualización, baja lógica, correo
duplicado, contraseña obligatoria, no autoescalada de rol propio,
control de acceso por rol). - \textbf{Trazabilidad:} → HU-022 → CU-22 →
\texttt{UsuarioController.\{listar,buscar,crear,actualizar,eliminar\}} →
\texttt{UsuarioService} → \texttt{GET\ /api/usuarios},
\texttt{GET\ /api/usuarios/\{id\}}, \texttt{POST\ /api/usuarios},
\texttt{PUT\ /api/usuarios/\{id\}},
\texttt{DELETE\ /api/usuarios/\{id\}}. - \textbf{Estado:} verificado.

\textbf{REQ-F-024 --- Gestión de vacunas} - \textbf{Tipo:} Funcional ·
\textbf{Prioridad:} Should - \textbf{Enunciado:} El sistema deberá
permitir registrar, consultar (de forma global, filtrada por mascota, o
por identificador), actualizar y dar de baja lógica las vacunas
aplicadas a una mascota, restringiendo la consulta de un usuario con rol
\texttt{ROLE\_DUENO} a las vacunas de sus propias mascotas. -
\textbf{Rationale:} nuevo requisito de la actividad de Unidad IV;
formaliza el módulo \texttt{VacunaController}/\texttt{VacunaService},
sin requisito formal previo en el SRS. - \textbf{Verificación:} Test
\texttt{VacunaControllerTest} (10 pruebas); colección Postman
\texttt{docs/postman/BIOPET-Vacunas.postman\_collection.json} (12
requests, ejecutable con Newman). La operación
\texttt{GET\ /api/vacunas/mascota/\{mascotaId\}} no tiene prueba
automatizada dedicada en \texttt{VacunaControllerTest}; las demás
operaciones sí. - \textbf{Trazabilidad:} → HU-023 → CU-23 →
\texttt{VacunaController.\{listar,listarPorMascota,buscar,crear,actualizar,eliminar\}}
→ \texttt{VacunaService} → \texttt{GET\ /api/vacunas},
\texttt{GET\ /api/vacunas/mascota/\{mascotaId\}},
\texttt{GET\ /api/vacunas/\{id\}}, \texttt{POST\ /api/vacunas},
\texttt{PUT\ /api/vacunas/\{id\}}, \texttt{DELETE\ /api/vacunas/\{id\}}.
- \textbf{Estado:} verificado.

\textbf{REQ-F-025 --- Consulta de información externa de especies} -
\textbf{Tipo:} Funcional · \textbf{Prioridad:} Could -
\textbf{Enunciado:} El sistema deberá permitir consultar información de
una especie (nombre científico, hábitat, dieta) desde un servicio
externo, utilizando una caché Redis para evitar llamadas repetidas al
servicio externo dentro de una ventana de tiempo configurable. -
\textbf{Rationale:} nuevo requisito de la actividad de Unidad IV;
integración con un proveedor externo (API Ninjas) vía
\texttt{ExternalApiClient}, con patrón \emph{cache-aside} en
\texttt{ExternalApiService} (TTL configurable mediante
\texttt{app.external-api.cache-ttl-seconds}, valor por defecto 600 s). -
\textbf{Verificación:} Test \texttt{ExternalApiServiceTest} (6 pruebas,
\texttt{Backend/src/test/java/com/biopet/integration/}) --- corrige una
afirmación de una versión anterior de este documento, que declaraba
``sin prueba automatizada dedicada''; esa clase de test sí existe en el
código actual. Evidencia empírica adicional en
\texttt{docs/u4/evidencias/fred/redis-cache-comparacion.md} (comparación
real de tiempos de respuesta con caché fría vs.~caché caliente). -
\textbf{Trazabilidad:} → HU-024 → CU-24 →
\texttt{ExternalApiController.infoEspecie} →
\texttt{ExternalApiService.obtenerInfoEspecie} →
\texttt{ExternalApiClient} → \texttt{GET\ /api/externa/especies}. -
\textbf{Estado:} verificado.

\subsubsection{3.2. Requisitos no funcionales
(REQ-NF)}\label{requisitos-no-funcionales-req-nf}

\textbf{REQ-NF-001 --- Rendimiento del listado de mascotas} -
\textbf{Categoría:} Rendimiento · \textbf{Prioridad:} Must -
\textbf{Enunciado:} El tiempo de respuesta del listado de mascotas no
deberá superar 200 ms (p95) con caché caliente ni 500 ms (p95) con caché
fría. - \textbf{Rationale:} heredado de RNF-01/RNF-WEB-04 de la Entrega
1A, con umbrales cuantitativos añadidos por el bloque C.1 de la Guía. -
\textbf{Verificación:} k6, 5 corridas en caliente y 5 en frío,
\texttt{docs/mediciones/perf/} (archivos
\texttt{k6-20260817T*-local-tls-v0.9.0-rc-*.json} y \texttt{REPORT.md}).
- \textbf{Estado:} verificado. \texttt{docs/mediciones/perf/REPORT.md}
registra p95 entre 9.7 ms y 22.13 ms según la corrida, muy por debajo de
los umbrales de 200 ms (caché caliente) y 500 ms (caché fría) exigidos
por este requisito, con 0.0 \% de error en todas las corridas.

\textbf{REQ-NF-002 --- Comunicación cifrada obligatoria} -
\textbf{Categoría:} Seguridad · \textbf{Prioridad:} Must -
\textbf{Enunciado:} Toda comunicación entre cliente y servidor deberá
realizarse mediante HTTPS/TLS. - \textbf{Rationale:} heredado de
RNF-02/RNF-WEB-02. - \textbf{Verificación:} \texttt{curl\ -v} mostrando
TLSv1.3 y suite AEAD; \texttt{SecurityHeadersTest}. -
\textbf{Trazabilidad:} \texttt{docker-compose.tls.yml},
\texttt{application-tls.yml}, \texttt{scripts/generate-dev-keystore.sh}.
- \textbf{Estado:} verificado en configuración TLS de desarrollo.

\textbf{REQ-NF-003 --- Expiración configurable de tokens JWT} -
\textbf{Categoría:} Seguridad · \textbf{Prioridad:} Must -
\textbf{Enunciado:} El sistema deberá emitir access tokens con
expiración de 1 hora y refresh tokens con expiración de 7 días, ambos
configurables mediante variables de entorno
(\texttt{JWT\_EXPIRATION\_MS}, \texttt{JWT\_REFRESH\_EXPIRATION\_MS}),
sin valores fijos en el código. - \textbf{Rationale:} heredado de
RNF-03/RNF-WEB-03. - \textbf{Verificación:} \texttt{JwtServiceTest};
inspección de \texttt{application.yml}. - \textbf{Estado:} verificado.

\textbf{REQ-NF-004 --- Claims estándar del JWT} - \textbf{Categoría:}
Seguridad · \textbf{Prioridad:} Must - \textbf{Enunciado:} Cada JWT
emitido deberá incluir los siete claims estándar (\texttt{iss},
\texttt{sub}, \texttt{aud}, \texttt{exp}, \texttt{nbf}, \texttt{iat},
\texttt{jti}) conforme al RFC 7519. - \textbf{Rationale:} refinamiento
de RNF-03 exigido por el bloque A.1 de la Guía; verificado directamente
en \texttt{JwtService.java} (issuer, audience configurables vía
\texttt{JWT\_ISSUER}/\texttt{JWT\_AUDIENCE}). - \textbf{Verificación:}
\texttt{JwtServiceTest}. - \textbf{Estado:} verificado.

\textbf{REQ-NF-005 --- Interfaz responsiva} - \textbf{Categoría:}
Usabilidad · \textbf{Prioridad:} Must - \textbf{Enunciado:} La interfaz
deberá adaptarse correctamente a resoluciones entre 320px y 1440px
(móvil, tablet, escritorio). - \textbf{Rationale:} heredado de
RNF-04/RNF-WEB-01. - \textbf{Verificación:} Lighthouse (perfil móvil
Slow 4G + perfil desktop) + inspección manual en Chrome DevTools a
320px, 768px y 1440px. - \textbf{Trazabilidad:}
\texttt{frontend/src/app/features/mascotas.component.ts}
(\texttt{.grid-mascotas} con \texttt{@media\ (max-width:\ 600px)}). -
\textbf{Estado:} verificado. \texttt{docs/mediciones/lighthouse/}
contiene 6 corridas originales (2026-08-01, perfil móvil) más 12
corridas de re-ejecución (2026-08-18, perfiles móvil y desktop,
\texttt{/login} y \texttt{/mascotas},
\texttt{lhci-20260818-0538-*.json}). Performance ≥90 en todos los casos
(100 en desktop, 90--94 en móvil), Accessibility 91 en todos los casos,
en ambos perfiles y ambas rutas (incluida la vista autenticada de
Mascotas).

\textbf{REQ-NF-006 --- Compatibilidad de navegadores} -
\textbf{Categoría:} Compatibilidad · \textbf{Prioridad:} Should -
\textbf{Enunciado:} La aplicación deberá funcionar correctamente en las
versiones modernas de Chrome, Firefox y Edge. - \textbf{Rationale:}
heredado de RNF-05/RNF-WEB-05. - \textbf{Verificación:} ejecución manual
de los flujos críticos en los tres navegadores. - \textbf{Estado:}
pendiente de evidencia archivada.

\textbf{REQ-NF-007 --- Disponibilidad durante evaluaciones académicas} -
\textbf{Categoría:} Disponibilidad · \textbf{Prioridad:} Must -
\textbf{Enunciado:} El sistema deberá estar operativo durante las
semanas de evaluación establecidas por la asignatura. -
\textbf{Rationale:} heredado de RNF-06. - \textbf{Verificación:}
demostración en vivo (\texttt{make\ up}) durante la semana de entrega;
en producción, healthcheck de Render (\texttt{docs/despliegue/}). -
\textbf{Estado:} \textbf{pendiente de confirmación explícita (Must \#1
de 2, ver sección 7)}. La afirmación ``verificado por diseño'' de una
versión anterior de este documento no está respaldada por un artefacto
concreto (log de uptime, captura de healthcheck con fecha, o corrida de
monitoreo real); es una inferencia razonable, no una verificación
empírica --- y coincide con lo que ya registra
\texttt{docs/trazabilidad/matriz.csv} para esta misma fila: estado
``implementado'' (no ``verificado''), con la nota explícita ``no existe
evidencia de disponibilidad sostenida durante una semana completa de
evaluación''. Se solicitó a Fred (responsable de despliegue/Render) el
nombre del artefacto o commit que la respalde antes de volver a marcarla
``verificado''.

\textbf{REQ-NF-008 --- Cifrado de contraseñas} - \textbf{Categoría:}
Seguridad · \textbf{Prioridad:} Must - \textbf{Enunciado:} Las
contraseñas de los usuarios deberán almacenarse cifradas mediante un
algoritmo hash seguro (BCrypt), nunca en texto plano. -
\textbf{Rationale:} heredado de RNF-07. - \textbf{Verificación:}
\texttt{AuthControllerTest}; inspección de \texttt{db/seed.sql} (hash
BCrypt del usuario admin) y de \texttt{SecurityConfig}
(\texttt{PasswordEncoder}). - \textbf{Estado:} verificado.

\textbf{REQ-NF-009 --- Registro de eventos de autenticación (OWASP A09)}
- \textbf{Categoría:} Seguridad · \textbf{Prioridad:} Must -
\textbf{Enunciado:} El sistema deberá registrar cada evento de
autenticación (login exitoso, login fallido, bloqueo por rate limit,
refresh, logout) con IP, marca de tiempo y sujeto (correo), sin exceder
200 caracteres por campo registrado. - \textbf{Rationale:} requisito
nuevo derivado del control OWASP A09, exigido por el bloque C.2 de la
Guía; ausente en el SRS original. Verificado directamente en
\texttt{AuthenticationAuditService.java}. - \textbf{Verificación:}
\texttt{AuthenticationAuditServiceTest}; captura de log real en
\texttt{docs/mediciones/sec/A09-logging.md}. - \textbf{Estado:}
verificado. \texttt{docs/mediciones/sec/A09-logging.md} incluye líneas
reales de log (\texttt{AUTH\_AUDIT\ event=LOGIN\_SUCCESS},
\texttt{LOGIN\_FAILURE}, \texttt{LOGIN\_RATE\_LIMITED}, con IP,
timestamp UTC y sujeto) y documenta la sanitización contra log forging.

\textbf{REQ-NF-010 --- Limitación de intentos de login (OWASP A07)} -
\textbf{Categoría:} Seguridad · \textbf{Prioridad:} Must -
\textbf{Enunciado:} El sistema deberá bloquear temporalmente los
intentos de login desde una misma IP tras 6 intentos fallidos
consecutivos dentro de una ventana de 15 minutos, respondiendo
\texttt{429\ Too\ Many\ Requests} durante el bloqueo, con parámetros
configurables mediante variables de entorno. - \textbf{Rationale:}
requisito nuevo derivado del control OWASP A07, exigido por el bloque
C.2 de la Guía; ausente en el SRS original. Verificado directamente en
\texttt{LoginRateLimiterService.java}
(\texttt{security.rate-limit.login.max-attempts}, \texttt{.window},
\texttt{.block-duration}). - \textbf{Verificación:}
\texttt{LoginRateLimiterServiceTest}; evidencia HTTP real en
\texttt{docs/mediciones/sec/A07-authentication.md}. - \textbf{Estado:}
verificado. \texttt{docs/mediciones/sec/A07-authentication.md} registra
los 5 primeros fallos respondiendo \texttt{401}, el 6º respondiendo
\texttt{429} con \texttt{ProblemDetail}
(\texttt{type=urn:biopet:error:rate-limited}) y cabecera
\texttt{Retry-After:\ 900} real, capturada con \texttt{curl}, no solo
con la prueba JUnit.

\textbf{REQ-NF-011 --- Arquitectura en capas} - \textbf{Categoría:}
Mantenibilidad · \textbf{Prioridad:} Must - \textbf{Enunciado:} La
arquitectura del backend deberá organizarse en capas separadas de
presentación (\texttt{controller}), lógica de negocio (\texttt{service})
y acceso a datos (\texttt{repository}/\texttt{entity}). -
\textbf{Rationale:} heredado de RNF-08. - \textbf{Verificación:}
revisión de la estructura de paquetes
\texttt{com.biopet.\{controller,service,repository,entity,dto\}}. -
\textbf{Estado:} verificado.

\textbf{REQ-NF-012 --- Caché del listado de mascotas con TTL externo} -
\textbf{Categoría:} Rendimiento · \textbf{Prioridad:} Should -
\textbf{Enunciado:} El listado paginado de mascotas deberá almacenarse
en caché Redis con un tiempo de vida (TTL) configurable mediante
variable de entorno (\texttt{CACHE\_TTL\_MS}, valor por defecto 300000
ms), sin valores fijos en código, y su tasa de aciertos (hit ratio)
deberá medirse y reportarse empíricamente. - \textbf{Rationale:}
heredado de la estrategia de caché ya definida en
\texttt{ADR-003-jwt-redis.md}; llevado a requisito no funcional
explícito exigido por el bloque A.1 de la Guía. - \textbf{Verificación:}
inspección de \texttt{application.yml}
(\texttt{spring.cache.redis.time-to-live}); evidencia cruda en
\texttt{docs/mediciones/redis/}. - \textbf{Estado:} verificado
parcialmente. \texttt{docs/mediciones/redis/} confirma el TTL real de la
clave de caché (195 s, mediante \texttt{TTL}) y su existencia bajo carga
(\texttt{DBSIZE}, \texttt{KEYS\ mascotas::*}). \textbf{No hay todavía
una medición explícita de hit ratio} (aciertos/total); solo se verificó
que la clave existe con el TTL configurado, no su tasa de aciertos a lo
largo del tiempo. Esto queda pendiente (ver Observaciones, sección 7).

\textbf{REQ-NF-013 --- Estrategia híbrida de acceso a datos (ORM +
procedimientos almacenados)} - \textbf{Categoría:} Mantenibilidad /
Seguridad · \textbf{Prioridad:} Must - \textbf{Enunciado:} Toda
operación de base de datos que no sea un CRUD elemental sobre una única
tabla (según la definición del bloque A.2.1 de la Guía) deberá
implementarse como función o procedimiento almacenado versionado en
\texttt{db/procs/}, invocado desde un repositorio Spring Data, quedando
prohibida la concatenación de entrada de usuario en cualquier fragmento
JPQL, HQL o SQL nativo. - \textbf{Rationale:} requisito nuevo exigido
explícitamente por el bloque A.2 de la Guía. Se marcó ``cumplida por
alcance actual'' en una versión anterior de este documento porque en ese
momento no existía ninguna operación no elemental; en v0.9.0-rc ya
existe una implementación real:
\texttt{fn\_resumen\_mascotas\_por\_especie} (ver REQ-F-021). -
\textbf{Verificación:} \texttt{ResumenEspeciesIntegrationTest},
\texttt{ProcedimientosBiopetIntegrationTest};
\texttt{scripts/audit-sql-dynamic.sh} (ausencia de SQL dinámico por
concatenación); \texttt{docs/basedatos/CATALOGO-SP.md}. -
\textbf{Estado:} verificado. La confirmación pendiente en una versión
anterior de este documento (que \texttt{scripts/audit-sql-dynamic.sh} se
hubiera ejecutado realmente sobre el estado v1.0.0 del código, con las
rutinas ya reclasificadas a \texttt{PROCEDURE}, y no solo sobre el
estado v0.9.0-rc anterior) ya se realizó: la reproducción se generó
\textbf{posteriormente}, sobre el commit exacto del tag histórico
\texttt{v1.0.0} (\texttt{0d5cd525ce648cca7219da204e16fa622e671a87}), en
un worktree aislado que no modificó el árbol de trabajo principal --- el
código verificado es el del tag; el log es evidencia generada después,
no un artefacto que existiera dentro del commit. Resultado:
\texttt{audit-sql-dynamic:\ 0\ hallazgos\ \ \ en\ 7\ archivo(s).\ PASA}
(código de salida 0). Evidencia archivada en
\texttt{docs/mediciones/sec/audit-sql-dynamic-v1.0.0.txt}
(\texttt{docs/mediciones/sec/audit-sql-dynamic-v1.0.0.sha256}).

\begin{center}\rule{0.5\linewidth}{0.5pt}\end{center}

\subsubsection{3.3. Criterios INVEST aplicados a las historias de
usuario}\label{criterios-invest-aplicados-a-las-historias-de-usuario}

Las 24 historias de usuario de
\texttt{docs/requisitos/historias/HistoriasUsuario.md} se evalúan aquí
explícitamente contra los seis criterios INVEST de Cohn (Independent,
Negotiable, Valuable, Estimable, Small, Testable), con justificación
breve y verificable contra campos que ya existen en cada historia (no se
introduce información nueva para esta evaluación).

\textbf{Negotiable, Valuable y Testable se cumplen en las 24 historias,
sin excepción}, verificado estructuralmente así:

\begin{itemize}
\tightlist
\item
  \textbf{Negotiable:} las 24 historias siguen el formato Connextra
  (``Como\ldots{} quiero\ldots{} para\ldots{}''), que expresa intención
  y valor, no una solución de diseño fija; el ``cómo'' (campos exactos,
  validaciones) queda abierto a negociación en el caso de uso y en la
  implementación.
\item
  \textbf{Valuable:} cada historia declara un rol beneficiario explícito
  (\texttt{Como\ \textless{}rol\textgreater{}}) y está trazada a un
  requisito funcional con \emph{rationale} propio en la sección 3.1,
  nunca a una tarea puramente técnica sin beneficiario.
\item
  \textbf{Testable:} las 24 historias incluyen un bloque
  \texttt{gherkin} con al menos un escenario \texttt{Given/When/Then}
  verificable; las 8 aún no implementadas (Should/Could pendientes) lo
  declaran explícitamente como ``comportamiento esperado, no
  implementado'', sin fingir que ya existe evidencia de ejecución.
\end{itemize}

\textbf{Independent, Estimable y Small varían por historia}
(justificación por fila, tomada del campo \emph{Dependencias} y del
estado real de cada historia):

\begin{longtable}[]{@{}
  >{\raggedright\arraybackslash}p{(\linewidth - 6\tabcolsep) * \real{0.2500}}
  >{\raggedright\arraybackslash}p{(\linewidth - 6\tabcolsep) * \real{0.2500}}
  >{\raggedright\arraybackslash}p{(\linewidth - 6\tabcolsep) * \real{0.2500}}
  >{\raggedright\arraybackslash}p{(\linewidth - 6\tabcolsep) * \real{0.2500}}@{}}
\toprule\noalign{}
\begin{minipage}[b]{\linewidth}\raggedright
HU
\end{minipage} & \begin{minipage}[b]{\linewidth}\raggedright
Independent
\end{minipage} & \begin{minipage}[b]{\linewidth}\raggedright
Estimable
\end{minipage} & \begin{minipage}[b]{\linewidth}\raggedright
Small
\end{minipage} \\
\midrule\noalign{}
\endhead
\bottomrule\noalign{}
\endlastfoot
HU-001 & Sí & Sí --- implementada, esfuerzo ya observado & Nota --- 2
REQ-F relacionados \\
HU-002 & Sí & Sí --- implementada, esfuerzo ya observado & Sí --- 1
REQ-F \\
HU-003 & Parcial --- depende de HU-002 (requiere haber iniciado sesión
previamente) & Sí --- implementada, esfuerzo ya observado & Sí --- 1
REQ-F \\
HU-004 & Parcial --- depende de HU-002 & Sí --- implementada, esfuerzo
ya observado & Sí --- 1 REQ-F \\
HU-005 & Parcial --- depende de HU-002 & Sí --- implementada, esfuerzo
ya observado & Sí --- 1 REQ-F \\
HU-006 & Parcial --- depende de HU-002 & Sí --- implementada, esfuerzo
ya observado & Sí --- 1 REQ-F \\
HU-007 & Parcial --- depende de HU-005 (control de acceso por rol) & Sí
--- implementada, esfuerzo ya observado & Sí --- 1 REQ-F \\
HU-008 & Parcial --- depende de HU-005 & Sí --- implementada, esfuerzo
ya observado & Sí --- 1 REQ-F \\
HU-009 & Parcial --- depende de HU-005, HU-008 & Sí --- implementada,
esfuerzo ya observado & Sí --- 1 REQ-F \\
HU-010 & Parcial --- depende de HU-005, HU-009 & Sí --- implementada,
esfuerzo ya observado & Sí --- 1 REQ-F \\
HU-011 & Parcial --- depende de HU-005, HU-009 & Sí --- implementada,
esfuerzo ya observado & Sí --- 1 REQ-F \\
HU-012 & Parcial --- depende de HU-007 (requiere que la mascota exista)
& Parcial --- hay código/tests reales para el subconjunto entregado (ver
Observaciones de la historia) & Sí --- 1 REQ-F \\
HU-013 & Parcial --- depende de HU-012 & Parcial --- estado
``Pendiente'', sin artefacto entregado que medir & Sí --- 1 REQ-F \\
HU-014 & Parcial --- depende de HU-007 & Parcial --- hay código/tests
reales para el subconjunto entregado (ver Observaciones de la historia)
& Sí --- 1 REQ-F \\
HU-015 & Parcial --- depende de HU-007 & Parcial --- estado
``Pendiente'', sin artefacto entregado que medir & Sí --- 1 REQ-F \\
HU-016 & Parcial --- depende de HU-012 (requiere historial clínico) &
Parcial --- estado ``Pendiente'', sin artefacto entregado que medir & Sí
--- 1 REQ-F \\
HU-017 & Parcial --- depende de HU-007 & Parcial --- estado
``Pendiente'', sin artefacto entregado que medir & Sí --- 1 REQ-F \\
HU-018 & Sí & Parcial --- estado ``Pendiente'', sin artefacto entregado
que medir & Sí --- 1 REQ-F \\
HU-019 & Parcial --- depende de HU-002 (requiere usuario autenticado) &
Parcial --- estado ``Parcial'', sin artefacto entregado que medir & Sí
--- 1 REQ-F \\
HU-020 & Parcial --- depende de HU-005 (control de acceso por rol),
HU-007 (requiere que existan mascotas registradas) & Sí ---
implementada, esfuerzo ya observado & Sí --- 1 REQ-F \\
HU-021 & Parcial --- depende de HU-001 (registro de usuario) & Parcial
--- estado ``Pendiente'', sin artefacto entregado que medir & Sí --- 1
REQ-F \\
HU-022 & Parcial --- depende de HU-005 (control de acceso por rol) & Sí
--- implementada, esfuerzo ya observado & Sí --- 1 REQ-F \\
HU-023 & Parcial --- depende de HU-007 (registro de mascota, requiere
que la mascota exista) & Sí --- implementada, esfuerzo ya observado & Sí
--- 1 REQ-F \\
HU-024 & Sí & Sí --- implementada, esfuerzo ya observado & Sí --- 1
REQ-F \\
\end{longtable}

\textbf{Lectura honesta de ``Independent''.} La mayoría de historias
depende de al menos otra (típicamente HU-002 ``sesión iniciada'' o
HU-005/HU-007 ``rol autorizado''/``mascota existente''), lo cual es
normal y esperado en un sistema con autenticación y CRUD relacional:
ninguna historia de gestión de un recurso es realmente independiente de
``existe sesión'' y ``el recurso padre existe''. Se documenta la
dependencia real en vez de declarar ``Independent: Sí'' de forma
genérica sin sustento, que habría sido la alternativa más fácil pero
menos honesta.

\textbf{Lectura honesta de ``Estimable''.} Las historias ya
\textbf{implementadas} (15 de 24) son estimables con certeza
retrospectiva: el código y los tests existen, su esfuerzo ya fue
observado. Las \textbf{parciales} (HU-012, HU-014: 2) son estimables
solo para el subconjunto ya entregado --- el esfuerzo del resto
(calendario interactivo, vista de historial consolidado) sigue siendo
una estimación no verificada. Las \textbf{pendientes} (7) no tienen
ningún artefacto que permita estimar con base empírica; su tamaño es una
suposición heredada de la Entrega 1A, no una medición.

\textbf{Nota sobre ``Small'' en HU-001.} Es la única historia que agrupa
dos REQ-F (registro exitoso + rechazo por correo duplicado) porque ambos
comparten el mismo flujo de entrada (\texttt{POST\ /api/auth/registro})
y serían artificialmente pequeños por separado; se mantiene como una
sola historia pequeña y cohesiva, no como una violación de INVEST.

\subsection{4. Trazabilidad (resumen)}\label{trazabilidad-resumen}

La matriz completa vive en \texttt{docs/trazabilidad/matriz.csv} (bloque
A.3.3 de la Guía), con la fila raíz por requisito y las columnas:
\texttt{id\_requisito}, \texttt{tipo}, \texttt{prioridad\_moscow},
\texttt{historia\_usuario}, \texttt{caso\_de\_uso},
\texttt{modulo\_codigo}, \texttt{endpoint\_api},
\texttt{prueba\_automatizada}, \texttt{tipo\_acceso},
\texttt{evidencia\_empirica}, \texttt{estado}. Este SRS es la fuente de
verdad para las columnas \texttt{id\_requisito}, \texttt{tipo},
\texttt{prioridad\_moscow} y \texttt{estado}; la matriz no debe declarar
un requisito que no exista aquí, ni omitir ninguno de los REQ-F/REQ-NF
listados en la sección 3.

\textbf{Correspondencia REQ-F ↔ HU ↔ CU (para llenar
\texttt{matriz.csv}):}

\begin{longtable}[]{@{}llllll@{}}
\toprule\noalign{}
REQ-F & HU & CU & REQ-F & HU & CU \\
\midrule\noalign{}
\endhead
\bottomrule\noalign{}
\endlastfoot
001 & HU-001 & CU-01 & 011 & HU-010 & CU-10 \\
002 & HU-001 & CU-01 & 012 & HU-011 & CU-11 \\
003 & HU-002 & CU-02 & 013 & HU-012 & CU-12 \\
004 & HU-003 & CU-03 & 014 & HU-013 & CU-13 \\
005 & HU-004 & CU-04 & 015 & HU-014 & CU-14 \\
006 & HU-005 & CU-05 & 016 & HU-015 & CU-15 \\
007 & HU-006 & CU-06 & 017 & HU-016 & CU-16 \\
008 & HU-007 & CU-07 & 018 & HU-017 & CU-17 \\
009 & HU-008 & CU-08 & 019 & HU-018 & CU-18 \\
010 & HU-009 & CU-09 & 020 & HU-019 & CU-19 \\
021 & HU-020 & CU-20 & 022 & HU-021 & CU-21 \\
\end{longtable}

\subsubsection{4.1. Trazabilidad histórica: identificadores originales →
identificadores actuales (cierre de
OBS-03)}\label{trazabilidad-histuxf3rica-identificadores-originales-identificadores-actuales-cierre-de-obs-03}

La retroalimentación oficial del SGA de la Entrega 1A señaló que ``los
RF-WEB se remapean a RF-16/RF-17 sin matriz de trazabilidad explícita en
esta entrega'' (ver \texttt{docs/observaciones/OBSERVACIONES.md},
OBS-03). Hasta esta revisión, el vínculo entre los identificadores
originales (\texttt{RF-NN} / \texttt{RF-WEB-NN}, Entrega 1A) y los
identificadores actuales (\texttt{REQ-F-NNN}) solo existía disperso en
el campo \emph{Rationale} de cada requisito individual (sección 3.1). La
siguiente tabla lo consolida en un único lugar explícito, sin omitir
ningún requisito funcional del sistema y sin inventar ningún origen que
no estuviera ya citado en este documento:

\begin{longtable}[]{@{}
  >{\raggedright\arraybackslash}p{(\linewidth - 10\tabcolsep) * \real{0.1667}}
  >{\raggedright\arraybackslash}p{(\linewidth - 10\tabcolsep) * \real{0.1667}}
  >{\raggedright\arraybackslash}p{(\linewidth - 10\tabcolsep) * \real{0.1667}}
  >{\raggedright\arraybackslash}p{(\linewidth - 10\tabcolsep) * \real{0.1667}}
  >{\raggedright\arraybackslash}p{(\linewidth - 10\tabcolsep) * \real{0.1667}}
  >{\raggedright\arraybackslash}p{(\linewidth - 10\tabcolsep) * \real{0.1667}}@{}}
\toprule\noalign{}
\begin{minipage}[b]{\linewidth}\raggedright
Identificador anterior (Entrega 1A)
\end{minipage} & \begin{minipage}[b]{\linewidth}\raggedright
Identificador actual
\end{minipage} & \begin{minipage}[b]{\linewidth}\raggedright
Descripción
\end{minipage} & \begin{minipage}[b]{\linewidth}\raggedright
Caso de uso
\end{minipage} & \begin{minipage}[b]{\linewidth}\raggedright
Historia
\end{minipage} & \begin{minipage}[b]{\linewidth}\raggedright
Estado
\end{minipage} \\
\midrule\noalign{}
\endhead
\bottomrule\noalign{}
\endlastfoot
RF-01 & REQ-F-001 & Registro de usuario dueño de mascota & CU-01 &
HU-001 & verificado \\
RF-01, RF-02 & REQ-F-008 & Creación de mascota asociada a un dueño
existente & CU-07 & HU-007 & verificado \\
RF-02 & REQ-F-009 & Listado paginado de mascotas activas por
propietario/rol & CU-08 & HU-008 & verificado \\
RF-02 & REQ-F-011 & Actualización de mascota con verificación de
propiedad & CU-10 & HU-010 & verificado \\
RF-03, RF-04 & REQ-F-013 & Registro y consulta de historial clínico &
CU-12 & HU-012 & pendiente \\
RF-05 & REQ-F-014 & Prescripción de medicamentos & CU-13 & HU-013 &
pendiente \\
RF-06 & REQ-F-015 & Gestión de citas veterinarias mediante calendario &
CU-14 & HU-014 & pendiente \\
\textbf{RF-07} & \textbf{REQ-F-022} & \textbf{Notificaciones al usuario
por correo electrónico (cierre de OBS-02)} & \textbf{CU-21} &
\textbf{HU-021} & \textbf{pendiente} \\
RF-08, RF-09 & REQ-F-016 & API de recepción de telemetría de
dispositivos IoT & CU-15 & HU-015 & pendiente \\
RF-10 & REQ-F-017 & Recomendaciones clínicas informativas (redacción
cerrada en OBS-04) & CU-16 & HU-016 & pendiente \\
RF-11, RF-12 & REQ-F-018 & Facturación digital y registro de pagos &
CU-17 & HU-017 & pendiente \\
RF-13, RF-WEB-02 & REQ-F-006 & Control de acceso por rol (RBAC) & CU-05
& HU-005 & verificado \\
RF-14 & REQ-F-019 & Reportes estadísticos exportables & CU-18 & HU-018 &
pendiente \\
RF-15 & REQ-F-020 & Auditoría de operaciones del sistema & CU-19 &
HU-019 & parcial \\
RF-16, RF-WEB-01 & REQ-F-003 & Autenticación mediante usuario y
contraseña & CU-02 & HU-002 & verificado \\
RF-17, RF-WEB-04 & REQ-F-005 & Cierre de sesión con revocación de token
& CU-04 & HU-004 & verificado \\
RNF-03, RNF-WEB-03 & REQ-F-004 & Renovación de sesión (refresh) & CU-03
& HU-003 & verificado \\
\end{longtable}

\textbf{Requisitos sin origen en la Entrega 1A} (nuevos, agregados
durante la Tercera Entrega para documentar funcionalidad ya presente en
el código, sin requisito previo que cerrar): \texttt{REQ-F-002} (rechazo
de correo duplicado), \texttt{REQ-F-007} (consulta del perfil propio),
\texttt{REQ-F-010} (consulta de mascota por id), \texttt{REQ-F-012}
(baja lógica de mascota) y \texttt{REQ-F-021} (resumen de mascotas por
especie, exigido por el bloque A.2.2 de la Guía de la Tercera Entrega).
Ninguno de estos sustituye ni duplica un identificador \texttt{RF-NN}
original.

Esta tabla es de solo lectura respecto a
\texttt{docs/trazabilidad/matriz.csv}: no lo reemplaza ni lo modifica
(ese archivo está fuera del alcance de los archivos autorizados para
este cierre de observaciones). Queda como acción de seguimiento
incorporar, en una futura revisión de \texttt{matriz.csv}, una columna
equivalente de origen histórico para \texttt{REQ-F-022}.

\subsection{5. Modelo de datos
(referencia)}\label{modelo-de-datos-referencia}

El diccionario de datos completo de las entidades implementadas
(\texttt{usuarios}, \texttt{mascotas}) está en
\texttt{docs/diccionario\_datos.md}, sincronizado con la migración
Flyway \texttt{database/migrations/V1\_\_schema\_inicial.sql}. El modelo
entidad-relación conceptual completo (incluyendo las entidades
pendientes: \texttt{Cita}, \texttt{Historial\_Clinico},
\texttt{Dispositivo\_IoT}, \texttt{Chat\_Triage},
\texttt{Producto\_Servicio}, \texttt{Factura},
\texttt{Detalle\_Factura}, \texttt{Proveedor}, \texttt{Mercaderia},
\texttt{Detalle\_Ingreso}) proviene de la Entrega 1A
(\texttt{PFC\_Entrega1A\_BMT.pdf}, sección 6) y se conserva como visión
de producto pendiente de materialización, sin duplicar aquí el DDL
completo.

\textbf{Diagrama entidad-relación (DER) --- dos artefactos distintos, no
intercambiables (cierre de OBS-05):}

\begin{itemize}
\tightlist
\item
  \texttt{docs/diagrams/der-biopet/der-biopet.png} --- renderizado
  generado a partir de la fuente Graphviz \texttt{der-biopet.dot}. Es un
  diagrama \textbf{dibujado}, no una exportación de una herramienta de
  modelado de base de datos.
\item
  \href{../observaciones/evidencias/DER-BIOPET-pgAdmin-ERD-Tool.png}{\texttt{docs/observaciones/evidencias/DER-BIOPET-pgAdmin-ERD-Tool.png}}
  --- exportación \textbf{real} generada desde \textbf{pgAdmin 4,
  herramienta ERD Tool}, solicitada explícitamente por la
  retroalimentación oficial del SGA de la Entrega 1B: \emph{``Exportar
  el DER desde pgAdmin 4 (ERD Tool) como PNG de alta resolución para el
  informe final''} (ver \texttt{docs/observaciones/OBSERVACIONES.md},
  OBS-05). PNG válido (firma de archivo verificada), 768×883 px,
  generado directamente sobre el esquema real de PostgreSQL, no sobre la
  fuente \texttt{.dot}.
\end{itemize}

\subsection{6. Interfaces de usuario
(referencia)}\label{interfaces-de-usuario-referencia}

Los wireframes conceptuales (login unificado, dashboard por rol,
historial clínico, gestión de citas) están documentados en la Entrega 1A
(\texttt{PFC\_Entrega1A\_BMT.pdf}, sección 7). La interfaz realmente
implementada en v1.0.0 está en \texttt{frontend/src/app/features/}:
\texttt{login.component.ts} (autenticación),
\texttt{mascotas.component.ts} (CRUD con paginación, formularios,
confirmación de borrado, resumen por especie y accesibilidad) y, nuevo
en Unidad IV, \texttt{vacunas.component.ts} (CRUD de vacunas).
\textbf{Citas, consultas, gestión administrativa de usuarios y consulta
externa de especies existen como API real (ver 3.1) pero sin pantalla
propia todavía} --- no hay componentes correspondientes en
\texttt{frontend/src/app/features/}. Ninguna de estas cinco pantallas
cuenta con wireframe formal actualizado --- se deja como observación
abierta en la sección 7.

\subsection{7. Observaciones e información
pendiente}\label{observaciones-e-informaciuxf3n-pendiente}

Esta sección se revisó contra el estado real del repositorio (v1.0.0).
Se mantiene el mismo principio que en versiones anteriores: no se marca
nada como resuelto sin evidencia verificable, y no se inventa contenido
para cerrar un pendiente.

\textbf{Resuelto desde la versión anterior de este documento:}

\begin{itemize}
\tightlist
\item
  \textbf{ADR de cambio de pila tecnológica} --- CERRADO.
  \texttt{docs/adr/ADR-002-pila-tecnologica.md} existe, está en estado
  ``Aceptado'' y documenta explícitamente la migración de ASP.NET Core 8
  (ADR-001, Entrega 1A) a Java 21 / Spring Boot 3.2, con contexto,
  evidencia verificable en \texttt{Backend/pom.xml} y alternativas
  consideradas. El señalamiento que traía \texttt{CAMBIOS-SRS.md} ya no
  aplica.
\item
  \textbf{Evidencia empírica de rendimiento (REQ-NF-001)} --- CERRADO.
  10 corridas k6 reales (5 caliente, 5 frío) en
  \texttt{docs/mediciones/perf/}, con p95 entre 9.7 ms y 22.13 ms, muy
  por debajo del umbral de 200/500 ms.
\item
  \textbf{Evidencia OWASP A07/A09 (REQ-NF-009, REQ-NF-010)} --- CERRADO.
  \texttt{docs/mediciones/sec/A07-authentication.md} y
  \texttt{A09-logging.md} contienen capturas HTTP (\texttt{curl}) y de
  log reales, no solo pruebas JUnit.
\item
  \textbf{Bitácora de observaciones}
  (\texttt{docs/observaciones/OBSERVACIONES.md}): existe, cubre OBS-01 a
  OBS-15 con fuente primaria (capturas SGA) y verificación cruzada
  contra \texttt{git\ log}. OBS-13 y OBS-14 (CRediT individual y
  afirmación falsa sobre ausencia de historial Git en
  \texttt{CONTRIBUTORS.md}) quedaron cerradas al corregirse
  \texttt{CONTRIBUTORS.md} con acuerdo del equipo completo.
\end{itemize}

\textbf{Cierre 2026-08-18 --- actualización tras la re-ejecución real de
Lighthouse.} Contrario a lo previsto en el cierre del 2026-08-17 (que
asumía que Fred no completaría sus pendientes), el equipo sí ejecutó la
re-corrida de Lighthouse y Fred sí entregó su bloque de trabajos
relacionados. Se actualiza el estado real:

\begin{itemize}
\tightlist
\item
  \textbf{Lighthouse --- perfil de escritorio y re-ejecución tras el fix
  de SEO --- CERRADO.} Re-corrida real del 2026-08-18
  (\texttt{docs/mediciones/lighthouse/lhci-20260818-0538-*.json}, 12
  archivos: perfil móvil y desktop, \texttt{/login} y
  \texttt{/mascotas}, 3 corridas cada uno). \textbf{SEO pasó de 82/100 a
  100/100 en las 12 corridas}, confirmando que los dos fixes aplicados
  (\texttt{meta\ description} en \texttt{index.html},
  \texttt{robots.txt} + \texttt{angular.json}) resolvían la causa raíz
  por completo. Performance ≥90 en todos los casos (100 en desktop,
  90--94 en móvil); Accessibility 91 en todos. REQ-NF-005 queda
  verificado en ambos perfiles.
\item
  \textbf{PRISMA (trabajos relacionados) --- CERRADO con 10 estudios.}
  Fred entregó su bloque (F20-F21) el 2026-08-18: 3 estudios incluidos
  (\texttt{docs/investigacion/handoff-fred-trabajos-relacionados.md}) +
  11 referencias BibTeX con DOI verificado
  (\texttt{handoff-fred-referencias.bib}). Total: 3 (Zaida) + 4 (Jaime)
  + 3 (Fred) = \textbf{10 estudios, meta de ≥8-9 cumplida}. Las 24
  referencias nuevas de los tres bloques ya están integradas en
  \texttt{docs/informe/referencias.bib} (39 entradas totales).
\end{itemize}

\textbf{Cierre 2026-09-03 --- revisión a v1.0.0 (recalificación).}

\begin{itemize}
\tightlist
\item
  \textbf{Versión y firma.} Este documento pasa de v0.9.0-rc a
  \textbf{v1.0.0}. El PDF \texttt{SRS-v1.0.0.pdf} \textbf{no está
  firmado} por el docente-director a la fecha de este cierre; se envía
  para aprobación y la firma se incorporará, con fecha real, solo cuando
  se reciba.
\item
  \textbf{INVEST --- CERRADO.} Se agrega la sección 3.3, con los seis
  criterios aplicados explícitamente a las 24 historias de usuario
  (justificación por criterio, tabla verificable fila por fila).
\item
  \textbf{Alineación de HU-012/HU-014 y REQ-F-013/REQ-F-015 con el
  código actual --- CERRADO.} Revisión contra el repositorio real
  detectó que ambos requisitos seguían marcados ``pendiente'' pese a que
  la Unidad IV ya entregó \texttt{CitaController}/\texttt{CitaService} y
  \texttt{ConsultaController}/\texttt{ConsultaService} con CRUD completo
  y tests dedicados. Se corrigen a ``parcial'' (no ``verificado'',
  porque el calendario interactivo y la vista de historial consolidado
  siguen sin implementar) en el SRS, en \texttt{HistoriasUsuario.md} y
  en \texttt{CasosDeUso.md}.
\item
  \textbf{REQ-F-025 corregido --- CERRADO.} Una versión anterior de este
  documento afirmaba que
  \texttt{ExternalApiController}/\texttt{ExternalApiService} no tenían
  prueba automatizada; sí la tienen (\texttt{ExternalApiServiceTest}, 6
  pruebas). Se corrige el estado a ``verificado''.
\item
  \textbf{Conteo de Must --- NO se declara 22/22.} De los 22 requisitos
  Must, \textbf{21 están verificados con evidencia concreta} (test
  automatizado, captura HTTP real, o corrida de medición archivada).
  \textbf{1 queda explícitamente ``implementado'' pero no verificado}
  (REQ-NF-007 --- disponibilidad durante evaluaciones académicas: el
  mecanismo de healthcheck funciona y está probado, pero no existe
  evidencia de disponibilidad sostenida durante una semana completa de
  evaluación). REQ-NF-013 (estrategia híbrida de acceso a datos), que
  también estuvo pendiente, ya se verificó:
  \texttt{scripts/audit-sql-dynamic.sh} se reprodujo sobre el commit
  exacto del tag histórico \texttt{v1.0.0}
  (\texttt{0d5cd525ce648cca7219da204e16fa622e671a87}) con resultado ``0
  hallazgos en 7 archivos. PASA'' (evidencia archivada en
  \texttt{docs/mediciones/sec/audit-sql-dynamic-v1.0.0.txt}). No se
  marca ``verificado'' sin ese respaldo, aunque una versión anterior de
  este documento sí los daba por verificados sin evidencia suficiente.
\end{itemize}

\textbf{Sigue como limitación declarada, no bloqueante para el cierre de
esta entrega:}

\begin{itemize}
\tightlist
\item
  \textbf{Hit ratio de caché Redis (REQ-NF-012).} Sigue sin una medición
  de \texttt{keyspace\_hits}/\texttt{keyspace\_misses}; solo hay
  evidencia de TTL y existencia de clave bajo carga.
\item
  \textbf{REQ-NF-006 (compatibilidad de navegadores).} Sin evidencia
  archivada de ejecución manual en Chrome/Firefox/Edge.
\item
  \textbf{Wireframes de la pantalla de Mascotas y de los módulos
  pendientes} (historial clínico, citas, facturación): siguen siendo los
  de la Entrega 1A. El producto real y su documentación de arquitectura
  (C4, DER, SRS) sí reflejan el sistema implementado.
\item
  \textbf{Pantallas sin interfaz propia (API real, sin UI).} Citas
  (calendario interactivo), historial clínico consolidado, gestión
  administrativa de usuarios y consulta externa de especies: las cuatro
  tienen backend real y verificado, pero ninguna tiene componente en
  \texttt{frontend/src/app/features/} todavía.
\item
  \textbf{Aislamiento de datos incompleto en Consultas (REQ-F-013).}
  \texttt{ConsultaService.listar()} no filtra por propietario para
  \texttt{ROLE\_DUENO}, a diferencia de
  \texttt{MascotaService}/\texttt{VacunaService} (registrado en
  \texttt{docs/trazabilidad/matriz.csv}). Un dueño de mascota
  autenticado podría ver consultas de mascotas ajenas a través de este
  endpoint. Se documenta aquí para que no quede oculto; su corrección es
  responsabilidad del propietario del módulo (Unidad IV), fuera de la
  zona de archivos de esta revisión.
\end{itemize}

Ninguno de estos puntos se cierra con datos inventados. Quedan como
limitaciones explícitas de la Entrega Final, consistentes con el
criterio de todo este documento.

\end{document}
